% =====================================================================
% neurips_broader_impact.tex
% Included from paper2_v3_submission.tex via \input{}.
% Expanded Broader Impact / dual-use statement (~1 page).
% =====================================================================

\section*{Broader Impact}
\label{sec:broader-impact}

This paper establishes a \emph{structural limit} on a specific class
of LLM safety interventions (continuous, utility-preserving wrapper
defenses). It does not propose new attacks, release new attack
corpora, or provide adversaries with new techniques. We summarize the
impact, dual-use, and deployment considerations below in the form
encouraged by the NeurIPS 2025 Broader Impact guidance.

\paragraph{Responsible-disclosure stance.}
The artifact is a \emph{theoretical impossibility}, not a
vulnerability report. There is no affected vendor, no proof-of-concept
exploit, and no coordinated-disclosure timeline to negotiate: the
claim is that \emph{any} defense satisfying continuity and
utility-preservation on a connected prompt space must leave boundary
prompts fixed. Responsible disclosure in this context means (i)
framing the result as a design constraint rather than a recipe for
bypass, (ii) providing the Lean-verified proofs so defenders can
evaluate the claim on its own terms, and (iii) accompanying every
empirical demonstration with explicit quantitative bounds (such as
$K^{*} = G/\ell - 1$) that tell defenders what the achievable region
looks like.

\paragraph{Who benefits.}
The primary beneficiaries are \textbf{defenders and safety
researchers}: the trilemma gives them a principled way to allocate
effort, a sanity check on new-wrapper proposals (if a published
wrapper claims continuity, utility preservation, \emph{and}
completeness, one of the three claims must be
wrong), and concrete parameter targets for the defense-design knobs
they actually control (Lipschitz constant, threshold geometry,
choice of feature space). Secondary beneficiaries are
\textbf{auditors and policy teams} who need defensible upper bounds
on what wrapper-class controls can achieve.

\paragraph{Who does \emph{not} benefit.}
An adversary seeking to jailbreak a particular deployed model gains
nothing actionable from this paper. The theorems do not identify
specific prompts, do not parameterize a model family, and do not
provide a search procedure over inputs. Our empirical forced-collapse
demonstration (\Cref{tab:forced-collapse,tab:paraphrase-unsafe})
operates entirely on \textbf{prompts already contained in public
jailbreak datasets}; no new adversarial text is produced by this
work that was not already available to adversaries via prior
publications and public leaderboards.

\paragraph{Dual-use analysis.}
We considered three plausible misuse vectors:
\begin{itemize}[nosep,leftmargin=1.6em]
  \item \emph{``The paper implies wrappers are useless, so vendors
        will stop deploying them.''} The paper explicitly rules this
        out (\Cref{sec:implications}): wrappers remain the
        cheapest risk-reduction primitive, they just cannot be
        total. This is analogous to the existence of adversarial
        examples not implying that image classifiers should be
        removed from production.
  \item \emph{``The paper tells attackers where to look.''} The
        boundary-fixation phenomenon is a consequence of continuity
        on a connected space, not a localized exploit. No attacker
        gains a search advantage from knowing that a boundary exists
        in general.
  \item \emph{``The paper's quantitative bounds help adversaries
        tune.''} The quantitative bounds
        (Theorems~\ref{thm:persistent}--\ref{thm:cone} and
        $K^{*} = G/\ell - 1$) are defender-facing: they require the
        \emph{wrapper's} Lipschitz constant and the \emph{surface's}
        boundary rate, both of which are available only to the
        operator running the safety stack, not to an external
        attacker probing the deployed system.
\end{itemize}

\paragraph{Mitigations.}
To minimize misuse potential, this paper (i) does not propose or
evaluate any new attack; (ii) releases only \emph{wrapper-side}
quantitative bounds (Lipschitz constant, persistent-unsafe measure,
threshold gap) rather than attack improvements; (iii) anonymizes the
artifact URL to comply with double-blind review and prevent
author-linked attack-adjacent discoverability; and (iv) pairs every
empirical table with the theorem it exercises and the inequality it
is expected to satisfy, so that readers cannot mistake the
demonstrations for generic red-teaming results.

\paragraph{Deployment guidance for practitioners.}
For a team deploying an LLM-backed product, the intended takeaway is
a \emph{design-constraint tool}, not an attack toolkit. Concretely:
(a) if a wrapper is continuous and utility-preserving, treat its
near-threshold behavior as structurally unreliable and do not rely on
it alone; (b) invest the saved attention in layers outside the
wrapper class---training-time alignment, output-side filters,
human-in-the-loop review of boundary cases, and architectural
controls on the underlying model; (c) measure $\hat L, \hat\ell,
\hat G, \hat K$ on your own surface and use $K^{*} = \hat G/\hat\ell
- 1$ as a planning number when specifying SLAs for the wrapper
layer; (d) if regulatory or contractual requirements demand
\emph{complete} wrapper safety, know that this forces either
discontinuity (e.g., hard blocklists with a non-negligible refusal
rate on benign inputs) or utility sacrifice (e.g., canonicalizing
prompt content), and plan the product experience accordingly.

\paragraph{Scope of this statement.}
Our results do not address fairness, bias, privacy, or environmental
impact of LLMs more broadly; they address a specific class of safety
controls and a specific structural limit on that class. Claims about
adequacy of alternative defense mechanisms (RLHF, constitutional
training, discrete filtering, ensembles) are out of scope.
